\chapter{이미 형식화한 것: 조합적 평면성과 쉬운 방향}
\label{chap:easy}

이 장의 항목은 모두 Lean으로 검증했다. 어려운 방향(\ref{chap:tools}장부터)은 여기 있는 정의와
도구 위에 세운다.

\section{순열의 순환}

\begin{definition}[순환 수]
\label{def:orbitCount}
\lean{Kuratowski.orbitCount}
\leanok
유한 집합 위의 순열 \(f\)에 대해 \(\operatorname{orb}(f)\)는 \(f\)의 순환 수다. 고정점도 길이 1인
순환 하나로 센다. Lean에서는 \texttt{SameCycle} 동치 관계로 나눈 몫의 크기(\texttt{Nat.card})다.
\end{definition}

\begin{lemma}[호환은 순환 수를 \(\pm 1\) 바꾼다]
\label{lem:swap}
\lean{Kuratowski.orbitCount_mul_swap_of_sameCycle, Kuratowski.orbitCount_mul_swap_of_not_sameCycle}
\leanok
\uses{def:orbitCount}
\(a \ne b\)가 \(f\)의 같은 순환에 있으면 \(\operatorname{orb}(f \cdot (a\ b)) = \operatorname{orb}(f) + 1\)이다.
\(a\), \(b\)가 서로 다른 순환에 있으면 \(\operatorname{orb}(f \cdot (a\ b)) + 1 = \operatorname{orb}(f)\)이다.
\end{lemma}

\begin{proof}
\leanok
서로 다른 순환에 있으면 \(f \cdot (a\ b)\)에서 두 순환이 합쳐지고 나머지 순환은 그대로다. Lean에서는
\(f \cdot (a\ b)\)의 한 걸음이 \(f\)의 같은 순환 안에서 움직이거나 \(a\), \(b\)의 순환 사이를 오간다는
"구역" 보조정리로 순환들을 대응시킨다.

같은 순환에 있으면 \(g = f \cdot (a\ b)\)로 두고 \(g\)에서 \(a\), \(b\)가 서로 다른 순환에 있음을 보인다.
그렇지 않다면 \(f\)와 \(g\)의 같은 순환 관계가 일치하고 \(g\)의 \(a\)를 지나는 순환은
\(f\)의 \(a\)를 지나는 순환에 호환 \((a\ b)\)를 곱한 것이 된다. 둘 다 같은 받침의 순환이라 부호가 같아야
하는데 호환은 부호를 바꾸므로 모순이다. 이제 \(f = g \cdot (a\ b)\)에 앞의 경우를 적용한다.
\end{proof}

\begin{definition}[점 하나 잘라 내기]
\label{def:cut}
\lean{Kuratowski.cut}
\leanok
\(\operatorname{cut}(f, a) = f \cdot (f^{-1}(a)\ a)\). 순환 \(\cdots \to p \to a \to q \to \cdots\)를
\(\cdots \to p \to q \to \cdots\)와 고정점 \(a\)로 바꾼다.
\end{definition}

\begin{lemma}[잘라 내기의 순환 수]
\label{lem:orbitCount_cut}
\lean{Kuratowski.orbitCount_cut, Kuratowski.sameCycle_cut}
\leanok
\uses{def:cut, def:orbitCount}
\(f(a) = a\)이면 \(\operatorname{cut}(f, a)\)의 순환 수는 \(f\)와 같고, 아니면 1 많다.
\(a\)가 아닌 두 점이 \(f\)의 같은 순환에 있으면 \(\operatorname{cut}(f, a)\)에서도 같은 순환에 있다.
\end{lemma}

\begin{proof}
\leanok
\uses{lem:swap}
\(f(a) \ne a\)이면 \(f^{-1}(a) \ne a\)이고 두 점은 같은 순환에 있으므로 보조정리~\ref{lem:swap}를 쓴다.
\end{proof}

\section{회전 체계와 평면성}

Mathlib에는 그래프의 평면성 정의가 없다. 위상적 정의(\(\mathbb{R}^2\)로의 매장)를 쓰려면 조르당
곡선 정리가 필요한데 이것도 Mathlib에 없다. 그래서 위상 그래프 이론의 표준 조합적 정의를 쓴다
(Mohar--Thomassen, \emph{Graphs on Surfaces}, 3장).

\begin{definition}[회전 체계와 면]
\label{def:rotation}
\lean{Kuratowski.RotationSystem, Kuratowski.RotationSystem.facePerm, Kuratowski.RotationSystem.Face}
\leanok
\uses{def:orbitCount}
그래프 \(G\)의 \emph{회전 체계} \(R\)은 다트(방향 있는 변) 위의 순열 \(\rho\)로서, 다트의 시작점을
보존하고 같은 꼭짓점에서 나가는 다트들이 \(\rho\)의 한 순환을 이루는 것이다. 즉 각 꼭짓점 둘레의
이웃 순서다. \emph{면 순열}은 \(\varphi(d) = \rho(\bar d)\)이고 (\(\bar d\)는 \(d\)의 역방향 다트),
\emph{면}은 \(\varphi\)의 순환이다.
\end{definition}

\begin{definition}[오일러 지표]
\label{def:eulerChar}
\lean{Kuratowski.RotationSystem.numFaces, Kuratowski.RotationSystem.eulerChar}
\leanok
\uses{def:rotation}
면의 개수 \(F = \operatorname{orb}(\varphi) + I\)이다. 여기서 \(I\)는 고립점의 수이고, 고립점 하나를
면 하나로 센다 (꼭짓점 하나짜리 구면 매장). 오일러 지표는 \(\chi(R) = V - E + F\)이다. 개수는 모두
\texttt{Nat.card}와 \texttt{Set.ncard}로 세므로 정의가 \texttt{Fintype}이나 \texttt{DecidableRel}
인스턴스에 의존하지 않는다.
\end{definition}

\begin{definition}[평면 그래프]
\label{def:IsPlanar}
\lean{Kuratowski.IsPlanar}
\leanok
\uses{def:eulerChar}
\(G\)가 \emph{평면 그래프}라는 것은 \(\chi(R) = 2c\)인 회전 체계 \(R\)이 있다는 것이다. \(c\)는 연결
성분의 수다.
\end{definition}

정리~\ref{thm:euler}에 따라 모든 회전 체계에서 \(\chi(R) \le 2c\)이므로, 이 등식은 모든 성분이 종수
0인 곡면(구면)에 매장된다는 뜻이다. 회전 체계와 곡면 매장의 대응(Heffter--Edmonds--Ringel)은
위상수학의 정리이며 형식화하지 않았다. 대신 \texttt{scripts/check\_definition.py}로 같은 정의를
작은 그래프에서 전수 검사했다. 꼭짓점 5개 이하인 그래프 가운데 정의상 비평면은 \(K_5\) 하나뿐이고,
Petersen 그래프, 정팔면체, 정육면체 등에서도 알려진 결과와 일치한다.

\section{수술과 오일러 부등식}

\begin{lemma}[연결 성분, 변, 고립점 세기]
\label{lem:components}
\lean{Kuratowski.card_connectedComponent_of_edgeSet_eq_empty, Kuratowski.card_connectedComponent_deleteEdges_le, Kuratowski.card_connectedComponent_deleteEdges_eq, Kuratowski.ncard_isolated_add_one_le, Kuratowski.card_connectedComponent_subdivideEdge, Kuratowski.ncard_edgeSet_subdivideEdge, Kuratowski.ncard_isolated_subdivideEdge}
\leanok
\uses{def:subdivideEdge}
변 \(uv\)를 지우면 연결 성분은 많아야 하나 늘고, 지운 뒤에도 \(u\)에서 \(v\)로 갈 수 있으면 그대로다.
변이 없으면 \(c = V\)이고, 변이 있으면 \(I + 1 \le c\)이다. 변 \(xy\)를 고립점 \(w\)로 세분하면 성분은 하나 줄고, 변은
하나 늘고, 고립점은 하나 준다.
\end{lemma}

\begin{proof}
\leanok
도달 가능성에 대한 직접 계산이다 (\texttt{Kuratowski/Components.lean}).
\end{proof}

\begin{theorem}[변 지우기]
\label{thm:del}
\lean{Kuratowski.RotationSystem.del, Kuratowski.RotationSystem.numFaces_del, Kuratowski.RotationSystem.numFaces_del_of_sameCycle, Kuratowski.RotationSystem.numFaces_del_of_not_sameCycle, Kuratowski.RotationSystem.reachable_del_of_not_sameCycle}
\leanok
\uses{def:eulerChar}
다트 \(d\)의 변을 지운 그래프에는 회전 체계 \(R \setminus d\)가 있고, 그 면 순열은
\(\varphi \cdot (d\ \bar d)\)에서 \(d\), \(\bar d\)를 잘라 낸 것이다. 따라서 \(d\)와 \(\bar d\)가 같은 면에
있으면 \(F\)가 1 늘고, 다른 면에 있으면 1 준다. 후자의 경우 지운 변의 두 끝은 지운 뒤에도 이어져 있다.
\end{theorem}

\begin{proof}
\leanok
\uses{lem:swap, lem:orbitCount_cut}
회전 \(\rho\)에서 \(d\), \(\bar d\)를 잘라 낸 것이 새 회전이다. 면 순열을 계산하면 위의 꼴이 되고,
순환 수는 보조정리~\ref{lem:swap}와 \ref{lem:orbitCount_cut}로 센다. 다른 면에 있는 경우에는 \(d\)의 면을 따라
돌면 지운 변을 쓰지 않는 길이 생긴다.
\end{proof}

\begin{theorem}[오일러 부등식]
\label{thm:euler}
\lean{Kuratowski.RotationSystem.eulerChar_le}
\leanok
\uses{def:eulerChar}
모든 회전 체계에서 \(\chi(R) \le 2c\)이다.
\end{theorem}

\begin{proof}
\leanok
\uses{thm:del, lem:components}
변의 수에 대한 귀납법. 변이 없으면 \(\chi = V + V = 2c\)이다. 변 하나를 지우면, 같은 면인 경우
\(\chi\)는 2 늘고 성분은 많아야 하나 는다. 다른 면인 경우 \(\chi\)도 성분 수도 그대로다.
\end{proof}

\begin{theorem}[부분 그래프]
\label{thm:mono}
\lean{Kuratowski.IsPlanar.mono, Kuratowski.IsPlanar.deleteEdge}
\leanok
\uses{def:IsPlanar}
평면 그래프의 부분 그래프는 평면 그래프다.
\end{theorem}

\begin{proof}
\leanok
\uses{thm:del, thm:euler, lem:components}
변 하나를 지우는 경우만 보면 된다. 같은 면인 경우 \(\chi\)가 2 늘었는데 정리~\ref{thm:euler}에 따라
\(\chi \le 2c'\)이고 \(c' \le c + 1\)이므로 \(c' = c + 1\)이고 등식이 유지된다. 다른 면인 경우 둘 다 그대로다.
\end{proof}

\section{세분}

\begin{definition}[변 세분하기]
\label{def:subdivideEdge}
\lean{Kuratowski.subdivideEdge}
\leanok
\(H\)의 변 \(xy\)를 꼭짓점 \(w\)로 나눈 그래프는 \(H - xy + xw + wy\)이다.
\end{definition}

\begin{definition}[세분, 위상 소수]
\label{def:subdiv}
\lean{Kuratowski.IsSubdivision, Kuratowski.ContainsSubdivision}
\leanok
\uses{def:subdivideEdge}
\(K\)(꼭짓점 집합 \(W\))를 \(f : W \hookrightarrow V\)로 \(V\)에 놓은 \(K.\mathrm{map}\,f\)에서 시작해,
변 \(xy\)를 아직 쓰이지 않은 꼭짓점 \(w\)(\(w \notin \operatorname{range} f\), \(w\)는 고립점)로
나누는 일을 반복해 얻는 그래프가 \(K\)의 \emph{세분}이다 (귀납 정의). \(G\)가 \(K\)의 세분을 부분
그래프로 가지면 \(K\)는 \(G\)의 \emph{위상 소수}다 (Diestel, \emph{Graph Theory}, 1.7절).
\end{definition}

\begin{theorem}[세분 되돌리기]
\label{thm:smooth}
\lean{Kuratowski.IsPlanar.of_subdivideEdge}
\leanok
\uses{def:IsPlanar, def:subdivideEdge}
\(H\)의 변 \(xy\)를 고립점 \(w\)로 세분한 그래프가 평면이면 \(H\)도 평면이다.
\end{theorem}

\begin{proof}
\leanok
\uses{lem:orbitCount_cut, lem:components}
\(w\)의 두 다트를 면 순열에서 잘라 내면 \(x \to w \to y\)로 돌던 면이 \(x \to y\)로 바로 간다. 다트를
\(H\)의 다트로 옮기면 \(H\)의 회전 체계가 되고, 면 순열의 순환 수는 그대로다. 한편 \(w\)가 고립점이 되어
면(고립점 하나를 면 하나로 센다)과 성분이 하나씩 늘고 변은 하나 주므로 등식이 유지된다.
\end{proof}

\section{쉬운 방향}

\begin{definition}[\(K_5\)와 \(K_{3,3}\)]
\label{def:K5K33}
\lean{Kuratowski.K5, Kuratowski.K33}
\leanok
\(K_5\)는 \(\mathrm{Fin}\,5\) 위의 완전 그래프, \(K_{3,3}\)은 \(\mathrm{Fin}\,3 \oplus \mathrm{Fin}\,3\) 위의
완전 이분 그래프다.
\end{definition}

\begin{lemma}[변 개수 상한]
\label{lem:edge_bound}
\lean{Kuratowski.IsPlanar.edge_bound, Kuratowski.IsPlanar.edge_bound_of_isBipartite}
\leanok
\uses{def:IsPlanar}
변이 있고 차수 1인 꼭짓점이 없는 평면 그래프는 \(E + 6 + 3I \le 3V\)를 만족한다. 이분 그래프이면
\(E + 4 + 2I \le 2V\)이다.
\end{lemma}

\begin{proof}
\leanok
\uses{lem:components}
단순 그래프라 길이 1인 면이 없고, 차수 1인 꼭짓점이 없으니 길이 2인 면도 없다. 따라서
\(3\operatorname{orb}(\varphi) \le 2E\)이다. 이분 그래프이면 길이 3인 면도 없어
\(4\operatorname{orb}(\varphi) \le 2E\)이다. 이것을 \(\chi = 2c\)와 \(I + 1 \le c\)에 넣는다.
\end{proof}

\begin{theorem}[\(K_5\), \(K_{3,3}\)은 평면이 아니다]
\label{thm:K5K33_nonplanar}
\lean{Kuratowski.K5_map_not_isPlanar, Kuratowski.K33_map_not_isPlanar, Kuratowski.K5_not_isPlanar, Kuratowski.K33_not_isPlanar}
\leanok
\uses{def:IsPlanar, def:K5K33}
어떤 \(f\)에 대해서도 \(K_5.\mathrm{map}\,f\)와 \(K_{3,3}.\mathrm{map}\,f\)는 평면이 아니다.
\end{theorem}

\begin{proof}
\leanok
\uses{lem:edge_bound}
\(K_5.\mathrm{map}\,f\)에서 \(E = 10\), \(I = V - 5\)이므로 \(10 + 6 + 3(V - 5) \le 3V\)가 거짓이다.
\(K_{3,3}.\mathrm{map}\,f\)에서는 \(E = 9\), \(I = V - 6\)이므로 \(9 + 4 + 2(V - 6) \le 2V\)가 거짓이다.
\end{proof}

\begin{theorem}[쿠라토프스키 정리, 쉬운 방향]
\label{thm:easy}
\lean{Kuratowski.not_isPlanar_of_containsSubdivision, Kuratowski.IsPlanar.not_containsSubdivision}
\leanok
\uses{def:IsPlanar, def:subdiv, def:K5K33}
\(K_5\)나 \(K_{3,3}\)의 세분을 부분 그래프로 갖는 유한 그래프는 평면 그래프가 아니다.
\end{theorem}

\begin{proof}
\leanok
\uses{thm:mono, thm:smooth, thm:K5K33_nonplanar}
\(G\)가 평면이면 세분 \(H \le G\)도 평면이고 (정리~\ref{thm:mono}), 세분을 한 단계씩 되돌린
\(K.\mathrm{map}\,f\)도 평면이다 (정리~\ref{thm:smooth}). 이것은 정리~\ref{thm:K5K33_nonplanar}에 어긋난다.
\end{proof}

\section{그 밖의 확인}

\begin{theorem}[\(K_4\)는 평면]
\label{thm:K4}
\lean{Kuratowski.K4_isPlanar}
\leanok
\uses{def:IsPlanar}
\(K_4\)는 평면 그래프다. 정의가 공허하지 않다는 확인이다.
\end{theorem}

\begin{proof}
\leanok
회전 순서를 표로 준다. 면 순열의 각 순환에 이름표를 붙여 순환이 4개임을 \texttt{decide}로 확인하면
\(4 - 6 + 4 = 2\)이다.
\end{proof}

\begin{theorem}[동형 불변]
\label{thm:iso}
\lean{Kuratowski.IsPlanar.of_iso}
\leanok
\uses{def:IsPlanar}
평면성은 그래프 동형에 대해 불변이다.
\end{theorem}

\begin{proof}
\leanok
동형은 다트, 변, 고립점, 연결 성분을 일대일로 옮기고, 옮긴 면 순열은 원래 면 순열의 켤레라 순환 수가
같다 (\texttt{orbitCount\_permCongr}).
\end{proof}
